% English translation of checked_pages/esq_119_059.tex; source: cote 119, site 59, PDF page 60, typescript page 57.
% AI-translated by GPT-6 Astra (OpenAI, ChatGPT Pro mode).
\scanpage{119/59}
\mspage{57}

In the case of stratified tame spaces properly so called (which
are not multiplicities strictly speaking), one can express
the ``unfolding'' of this structure conveniently, i.e.
the system of spaces $X_{i_o,\ldots,i_n}$, in terms of the tame
space $X_*$\frapparatus{V-FR-059-001} that is the sum of the $X_i$, which is endowed\frapparatus{V-FR-059-002} with the \uline{structure of an}
\uline{ordered object} (in the category of tame spaces), having as the graph
$X_{**}$ of the order relation the sum of the $X_{i,j}$ and the $X_i$
(the latter constituting the diagonal). Among the essential
properties of this ordered structure, let us note only
here that $pr_1 : X_{**} \longrightarrow X_*$ is a proper (locally trivial)
fibration, and $pr_2 : X_{**} \longrightarrow X_*$ is a ``bordering'' embedding.\frapparatus{V-FR-059-003}
There is an analogous interpretation of the unfolding of a stratified
tame multiplicity, in terms of a \uline{category structure}
(replacing a simple ordered structure) ``in the sense of tame
multiplicities'', whose composition map is given
by the morphisms $X_{ijk} \longrightarrow X_{ik}$ above.

\margin{59}
